\chapter{危害公共安全罪}

% Source PDF physical page 316; see audit/source_coverage.json.
\begin{EditionOpeningPanel}
1.知识体系。前两讲是侵犯个人法益的犯罪。自本讲至第二十二讲是侵犯社会法益的犯罪。社会法益包括：公共安全、经济秩序、社会秩序。这三讲罪名数量多，主抓重点罪名。

2.知识储备：危害公共安全的含义。

\InsertDiagram{figures/manual_table_p316_1.tex}
\end{EditionOpeningPanel}

\SourceFootnote{\textbf{表注①}\quad “公众”中，（1）“多数人”的含义，不是指特定的多数人（灭门案），而是指有危及多数人的可能性（食堂锅里投毒案）。（2）“不特定人”的含义，不是指危及对象的不特定（往楼下人群扔砖块），而是指危及范围的不特定，有危及多数人的可能性（往楼下人群扔手雷）。}

% Source PDF physical page 317; see audit/source_coverage.json.
\section{危险方法型犯罪}

\subsection*{一、放火罪}

第114条放火、决水、爆炸以及投放毒害性、放射性、传染病病原体等物质或者以其他危险方法危害公共安全，\CJKunderwave{尚未造成严重后果的}，处三年以上十年以下有期徒刑。

第115条（第1款）放火、决水、爆炸以及投放毒害性、放射性、传染病病原体等物质或者以其他危险方法\CJKunderwave{致人重伤、死亡或者使公私财产遭受重大损失的}，处十年以上有期徒刑、无期徒刑或者死刑。

（第2款）\CJKunderwave{过失犯前款罪的}，处三年以上七年以下有期徒刑；情节较轻的，处三年以下有期徒刑或者拘役。

\subsubsection*{（一）既未遂问题}

第114条规定的是基本犯。第115条第1款规定的是结果加重犯。第115条第2款规定的是过失犯罪，如失火罪、过失爆炸罪、过失投放危险物质罪、过失以危险方法危害公共安全罪。

1.第114条基本犯的既未遂问题\footnote{按照正常原理，产生危险是故意犯罪的成立条件，造成实害结果才是既遂条件，所以，第114条是关于危险犯的规定，是第115条第1款（对加重结果持故意的结果加重犯）的未遂犯规定。既然第114条是未遂犯规定，就不应再有自己的既遂和未遂。

但是，第114条基本犯有自已的既遂和未遂。原因是，第一，立法者对该危险犯规定了独立的条文和独立的法定刑（也即第114条），该独立法定刑是针对既遂而言的。第二，结果加重犯的基本犯本来就拥有自己的既未遂问题。例如，强奸罪（故意）致人重伤中，基本犯（强奸罪）有自己的既未遂，加重犯“故意致人重伤”也有自己的既未遂。第114条是基本犯，应拥有自己的既未遂。

不过，该基本犯的既遂标准就不能是造成实害结果（伤亡、重大损失），因为该实害结果是第115条对加重结果持故意的结果加重犯的既遂标准。}

着手标准：实施足以点燃的行为。既遂标准：独立燃烧，并对公共安全制造了严重危险。独立燃烧，是指不需借助外界条件，火势已经开始独立燃烧。

2.第115条第1款结果加重犯的既未遂问题

该结果加重犯包括两种情形。（1）对加重结果持过失心理，比如放火罪（过失）致人死亡。由于“过失致人死亡”只有具备与否问题，所以不存在既未遂问题。

（2）对加重结果持故意心理，比如放火罪（故意）致人死亡。这种结果加重犯有既未遂问题，也即造成严重后果（伤亡、重大损失）才既遂。若未造成严重后果，构成未遂。\footnote{但对该未遂，不用根据总则未遂犯的规定进行从宽处罚（可以从轻或减轻处罚），应直接适用第114条规定即可（该规定对行为人更有利，因为是必然减轻处罚）。就此而言，第114条是第115条第1款的未遂犯的特别规定。}

提示：爆炸罪，投放危险物质罪，以危险方法危害公共安全罪的既未遂问题，也需要根据第114条（基本犯）和第115条第1款（结果加重犯）分别分析。

\subsubsection*{（二）罪数问题}

四罪（放火罪，爆炸罪，投放危险物质罪，以危险方法危害公共安全罪）与故意杀人罪的关系（2023年试题）

% Source PDF physical page 318; see audit/source_coverage.json.
\InsertDiagram{figures/manual_table_p318_1.tex}
\SourceFootnote{\textbf{表注①}\quad 实际处罚方案的选择问题。先回顾下想象竞合犯与结果加重犯的转化过程。例如，鲁智深拳打镇关西，一个行为同时触犯故意伤害罪和过失致人死亡罪，想象竞合，择一重罪论处，定故意伤害罪。但是，立法者认为该处罚结论过轻，进行特殊规定，也即将过失致人死亡作为故意伤害罪的加重处罚结果，生成结果加重犯“故意伤害罪致人死亡”，该处罚绝对比单纯的故意伤害罪重。

同理，甲的一个放火行为同时触犯放火罪和故意杀人罪，想象竞合，择一重罪论处，定故意杀人罪。但是，立法者认为该处罚结论较轻，进行特殊规定，将故意杀人作为放火罪的加重处罚结果，生成结果加重犯“放火罪（故意）致人死亡”。

问题在于，立法者给该结果加重犯配置的刑罚并不必然比故意杀人罪更重，也即立法者的配套措施没做好。由此导致，在实际处罚方案的选择上需要根据具体案情情节具体判断。如果放火罪（故意）致人死亡的处罚比故意杀人罪重，就按照放火罪（故意）致人死亡处理。如果放火罪（故意）致人死亡的处罚比故意杀人罪轻，就按照故意杀人罪处理，也即放火罪和故意杀人罪想象竞合，择一重罪论处，定故意杀人罪，此时不需要转化为结果加重犯。

提示，在抢劫罪（故意）致人死亡中，也存在上述问题，分析原理相同。参见财产犯罪一讲中抢劫罪一节的脚注。}

\subsection*{二、投放危险物质罪}

投放危险物质罪，是指故意投放毒害性、放射性、传染病病原体等危险物质，危害公共安全的行为。（2017年第57题）

1.危险物质，不仅包括毒害性物质，还包括放射性、传染病病原体等物质，也包括其他能危害公共安全的有害物质。本罪原为投毒罪，现为投放危险物质罪。例如，甲将患有狂犬病的狗扔到小区的垃圾站，属于投放危险物质。

2.罪名比较。投放虚假危险物质罪（第291条之一）：投放虚假的爆炸性、毒害性、放射性、传染病病原体等物质，严重扰乱社会秩序的行为。该罪不是危害公共安全罪，而是扰乱社会管理秩序罪。

% Source PDF physical page 319; see audit/source_coverage.json.
\subsection*{三、以危险方法危害公共安全罪}

1.构成要件

\InsertDiagram{figures/manual_table_p319_1.tex}

2.根据司法解释，下列情形可以定以危险方法危害公共安全罪：

（1）在多人通行的场所私拉电网，危害公共安全。

（2）在火灾现场破坏消防器材，危害公共安全。

（3）破坏矿井通风设备，危害公共安全。

（4）在具有瓦斯爆炸高度危险的情形下，下令多人下井采煤。

% Source PDF physical page 320; see audit/source_coverage.json.
（5）危险驾驶。第一，在高速公路上逆向高速行驶。第二，为制造社会影响，驾驶机动车在公共场所任意冲撞，危害公共安全。第三，醉酒驾车，肇事后继续驾车冲撞，放任危害后果的发生，造成重大伤亡。（2012年第15题）

\section{破坏型犯罪}

[问题］狗蛋故意砸坏行驶中的公交车的乘客座椅，是否构成破坏交通工具罪？（破坏座椅案)

\subsection*{一、破坏交通工具罪}

第116条破坏火车、汽车、电车、船只、航空器，足以使火车、汽车、电车、船只、航空器发生倾覆、毁坏危险，\CJKunderwave{尚未造成严重后果}的，处三年以上十年以下有期徒刑。

第119条破坏\CJKunderwave{交通工具、交通设施}、电力设备、燃气设备、易燃易爆设备，\CJKunderwave{造成严}\CJKunderwave{重后果}的，处十年以上有期徒刑、无期徒刑或者死刑。

过失犯前款罪的，处三年以上七年以下有期徒刑；情节较轻的，处三年以下有期徒刑或者拘役。

\subsubsection*{（一）交通工具的范围}

1.这里的“汽车”可作扩大解释，包括大型拖拉机。这里的“汽车”不仅包括公共汽车，也包括私家汽车。

2.马车、人力三轮车、自行车、电动自行车、摩托车，不属于本罪的交通工具。这是因为，第一，法条没有规定摩托车等；第二，破坏这些小型交通工具，不会严重危害公共安全。（2022年试题）

\subsubsection*{（二）成立条件与既遂标准}

1.第116条基本犯（破坏交通工具罪）

（1）成立条件：破坏行为对交通安全制造了危险。

例1，本节开头“破坏座椅案”中，行为没有对交通安全制造危险，不构成本罪。狗蛋如果破坏了刹车装置或方向盘，则制造了危险，构成本罪。狗蛋如果将车轮全部卸掉，汽车躺在地上，则不会对交通安全制造危险，不构成本罪，但可以构成故意毁坏财物罪。

例2（2017年第12题），甲乘坐飞机，飞机在等候滑出停机位的指令，空姐告知“不得打开安全门”，甲仍拧开安全门，导致飞机延误起飞。由于不足以发生倾覆危险，不构成本罪。如果飞机马上起飞，或已经起飞，则甲的行为构成本罪。

（2）着手标准：开始破坏交通工具，并对交通安全制造了危险。

（3）既遂标准：破坏了交通工具，导致交通工具丧失功能，并对交通安全制造了严重危险。

例1，甲正在破坏仇人乙的轿车的刹车装置，被保安发现，及时制止，甲未能破坏刹车装置。甲构成本罪的未遂。（2018年试题）

例2，甲的卡车因为车祸交付修理厂维修，修理工乙发现甲是自己的情敌，便在维修过程中故意破坏输油装置，导致该装置丧失功能，然后交付甲使用。乙构成本罪的既遂。（2006年第60题）

% Source PDF physical page 321; see audit/source_coverage.json.
2.第119条结果加重犯（破坏交通工具罪导致伤亡等严重后果）

该结果加重犯的主要情形是对加重结果持故意心理，所以存在既遂与未遂问题。

既遂标准：发生严重后果，例如车辆倾覆，酿成车祸，造成伤亡、重大损失等。

例如，电视剧《狂飙》里，老默成功破坏了村长李有田的轿车的刹车装置。假如李有田车技很好，竟然安全停靠，那么老默构成第119条结果加重犯的未遂。对此，直接适用第116条即可。第116条就是第119条的未遂规定。老默适用第116条，也构成第116条破坏交通工具罪的既遂。

\subsection*{二、破坏交通设施罪}

第117条破坏轨道、桥梁、隧道、公路、机场、航道、灯塔、标志或者进行其他破坏活动，足以使火车、汽车、电车、船只、航空器发生倾覆、毁坏危险，\CJKunderwave{尚未造成严重后}\CJKunderwave{果}的，处三年以上十年以下有期徒刑。

\subsubsection*{（一）行为方式}

1．“破坏”既包括物理性损毁，也包括使道路通行功能丧失。例如，在公路上撒钉子，放大石头，也属于破坏交通设施。（2018年试题）

2.盗窃机动车道路的井盖，同时触犯盗窃罪和破坏交通设施罪，想象竞合，择一重罪论处。\footnote{盗窃非机动车道、人行道的井盖，除了构成盗窃罪，司法解释（多数观点）认为还构成以危险方法危害公共安全罪；少数观点认为，不构成以危险方法危害公共安全罪，因为制造的危险没有危及多数人的可能性，没有危害公共安全。}

\subsubsection*{（二）既未遂问题}

1.第117条基本犯（破坏交通设施罪）

着手标准：开始破坏交通设施，并对交通安全制造了危险。

既遂标准：破坏了交通设施，导致交通设施丧失功能，并对交通安全制造了危险。

例如，甲正在破坏铁轨，被人发现，及时制止，未能破坏铁轨。甲构成本罪的未遂。

2.第119条结果加重犯（破坏交通设施罪导致伤亡等严重后果）

既遂标准：发生严重后果，例如伤亡结果、重大损失等。

例如，甲破坏了铁轨，铁轨断裂了一段，但是养路工及时发现，及时修复，未发生车祸。甲构成第119条结果加重犯的未遂。对此，直接适用第117条的既遂即可。

\subsection*{三、破坏电力设备罪（第118条）}

这是指，故意破坏电力设备，危害公共电力安全的行为。

行为对象：已经交付使用的电力设备，不要求通电。例如，某公司将电力设备安装完成，但未交付使用。甲为泄愤，损毁了电力设备。甲不构成本罪，构成故意毁坏财物罪。（2025年试题）

\section{交通型犯罪}

[问题］狗蛋驾车不慎撞死一人，重伤一人，狗蛋不知道，继续往前开。对狗蛋该如何处理？（不知撞人案）

% Source PDF physical page 322; see audit/source_coverage.json.
\subsection*{一、交通肇事罪}

第133条违反交通运输管理法规，因而发生重大事故，致人重伤、死亡或者使公私财产遭受重大损失的，处三年以下有期徒刑或者拘役；交通运输\CJKunderwave{肇事后逃逸}或者有其他特别恶劣情节的，处三年以上七年以下有期徒刑；\CJKunderwave{因逃逸致人死亡的}，处七年以上有期徒刑。

\subsubsection*{（一）}

成立条件（基础刑是3年以下有期徒刑）

\InsertDiagram{figures/manual_table_p322_1.tex}

\subsubsection*{（二）法定刑升格条件一：“肇事后逃逸”（升格刑是3年至7年有期徒刑)}

1.司法解释规定，“肇事后逃逸”，是指行为人为逃避法律追究而逃跑的行为。\footnote{2000年11月21日最高人民法院《关于审理交通肇事刑事案件具体应用法律若干问题的解释》。}

2.前提认识。行为人需要明知发生了交通事故，否则不能认定为逃逸。例如，本节开头“不知撞人案”中，狗蛋构成交通肇事罪，但不知道自已撞了人，还继续行驶，不构成“肇事后逃逸”。

\subsubsection*{（三）法定刑升格条件二：“因逃逸致人死亡”（升格刑是7年至15年有期徒刑）}

“因逃逸致人死亡”，是指为逃避法律追究而逃跑、不救人，导致人死亡，是一种不作

% Source PDF physical page 323; see audit/source_coverage.json.
为犯罪。对此，可以按照不作为犯罪的成立条件来分析。

1.作为义务（应为）

作为义务（救人义务）源于先行行为（肇事行为）。例如，甲开车不慎撞死乙、丙。甲误以为丙没有死亡，为逃避救助义务而逃逸。因为丙已死，不存在救人义务，甲不构成“因逃逸致人死亡”。

2.作为可能性（能为）

例如，甲开车不慎撞死乙、重伤丙。甲想救丙，但是丙的亲戚殴打甲，导致甲无法救助。甲驾车至交警大队报案。丙因未得到及时救助而死亡。甲不构成“因逃逸致人死亡”

3.不履行导致死亡结果（不为）

（1）结果避免可能性。例如，甲开车不慎撞死乙，重伤丙，严重程度是再怎么及时救助也救不活。甲不救丙，逃逸，丙因伤势过重而死亡。即使甲履行了救助义务，也不能避免丙的死亡，因此，丙的死亡不归属于甲的不作为，而归属于甲的前面肇事行为。甲不构成“因逃逸致人死亡”。（2013年不定项选择题）

（2）因果关系。最全面复杂的行为结构：肇事行为+不作为（不救人）+介入因素一死亡结果。

例1，甲开车不慎撞死乙，重伤丙，丙倒在路中央，甲逃逸。5分钟后，丁驾车不慎轧到丙，丙死亡。第一，死亡结果与丁的碾轧有因果关系。第二，介入因素（丁的不慎轧）不异常，是甲的行为引发的，因此丙的死亡结果与甲的行为（一是肇事行为，二是逃逸、不作为）有因果关系。汇总下来，死亡结果与丁的一个行为、甲的两个行为有因果关系，属于三因一果。第三，甲因死亡结果只能受一次处罚，如果处罚两次，则违反了禁止重复处罚原则。由于甲的肇事行为已经构成交通肇事罪，不需要将死亡结果用给肇事行为，所以用给不作为（逃逸），基于此，甲构成“因逃逸致人死亡”。第四，丁主观上若有过失，则构成交通肇事罪。

例2（无法查明），甲开车不慎撞倒乙，乙倒在路中央，甲逃离现场。10分钟后，丙开车不慎轧过去。20分钟后，救护人员发现乙死亡，但无法查明死亡时间。可能的死亡时间点有三个：一是甲的不慎撞倒行为，直接导致死亡。二是，不慎撞倒没撞死，逃逸期间，在丙轧之前，乙死亡。三是，丙轧之前没死，丙轧之后死亡。死亡时间无法查明时，死得越早对大家越有利。因此，认定为甲的不慎撞倒行为直接导致死亡。基于此分析，甲成立交通肇事罪，属于“肇事后逃逸”，但不属于“因逃逸致人死亡”。丙不成立交通肇事罪。

4.等价性

“因逃逸致人死亡”是不作为致人死亡。不作为的等价性有两个程度等级：第一，单纯逃逸，相当于遗弃罪程度，定“因逃逸致人死亡”。第二，隐藏伤者，排除他人救助可能性，达到了不作为杀人程度，定故意杀人罪。

例如，甲开车不慎撞死乙，重伤了丙，故意不救助丙，并且将丙抬到一个草丛中，丙因未得到及时救助而死亡。甲的前一行为构成交通肇事罪，后一行为中，甲的不作为在等价性上已经达到故意杀人罪的程度，因为自已不救，还排除了他人救助的可能性。因此对甲不再以“因逃逸致人死亡”论处，而定不作为的故意杀人罪，与交通肇事罪并罚。如果甲故意将丙扔到河里淹死，则构成作为的故意杀人罪，与交通肇事罪并罚。如果不处罚交通肇事罪，则遗漏评价了乙的死亡。（2007年第9题）

% Source PDF physical page 324; see audit/source_coverage.json.
5.主观要件。第一，要求知道自已肇事了。第二，对于逃逸不救助致人死亡，既可以是故意为之，也可以过失为之。

例1，甲开车不慎撞死乙，重伤了丙，明知不救丙，丙就会死，仍逃离现场，丙因未得到及时救助而死亡。甲构成交通肇事罪，并构成“因逃逸致人死亡”，是故意为之。

例2（2007年第9题），甲开车不慎撞死乙，重伤了丙，以为丙也死亡，便驾车离去，丙因未能得到及时救助而死亡。甲构成交通肇事罪，并构成“因逃逸致人死亡”，是过失为之。\footnote{关于“因逃逸致人死亡”的体系地位，存在观点展示。（1）主流观点认为（前提条件说），法条规定，“因逃逸致人死亡”是交通肇事罪的量刑情节（法定加重处罚情节，法定刑升格条件），因此，适用“因逃逸致人死亡”的前提是，前面肇事行为已经成立交通肇事罪。在此基础上，才可能适用“因逃逸致人死亡”。理由是，先有定罪，后才会有量刑；没有定罪，不能直接量刑。参见张明楷：《刑法学》（第六版），第928页。

（2）少数观点认为（非前提条件说），适用“因逃逸致人死亡”不要求前面肇事行为成立交通肇事罪。这种观点的问题在于，“因逃逸致人死亡”不是个罪名，只是个量刑情节。适用“因逃逸致人死亡”时，总得给被告人定个罪名。而这种观点又认为前面肇事行为不需要成立交通肇事罪。那么给被告人认定“因逃逸致人死亡”时，定什么罪名？有法官为非前提条件说提供的理由是，连环碾轧案中，根本无法判断、无法查明第一次碾轧行为是否导致重伤。由此，无法认定第一次碾轧行为成立交通肇事罪。所以，不能要求第一次碾轧行为成立交通肇事罪。（参见《刑事审判参考》第1118号，邵某交通肇事案）。

这种看法有误。即使不能认定第一次碾轧导致重伤，但照样可以给第一次碾轧行为定交通肇事罪。例如，甲驾车不慎撞倒丙，甲逃逸，五分钟后，乙驾车不慎轧过丙，丙才死亡。无法查明甲是否将丙撞成重伤。因此，无法按照“一个重伤+逃逸”给其定交通肇事罪。但是，最终死亡结果与甲的撞击行为有因果关系，因为介入因素“乙的碾轧”并不异常。所以，甲的撞击行为成立交通肇事罪。提示，成立交通肇事罪所要求的一个死亡结果，并不要求当场死亡，可以是介入因素之后才死亡。有些因果关系，因与果之间可以存在较长时间间隔。这是2013年卷二不定项题的官方观点。}

\subsubsection*{（四）共犯问题}

1.共同过失肇事

（1）监督过失。例如，甲看到乙醉酒，仍将自已的车辆借给乙驾驶，乙酿成车祸，致人死亡。乙构成交通肇事罪。甲对自己的车辆具有监督过失，构成交通肇事罪。

（2）指挥过失。例1，司机送领导去机场，领导指使司机闯红灯。司机照办，不慎撞死行人。司机构成交通肇事罪。领导有指挥过失，也构成交通肇事罪。例2，女友看到男友醉酒，仍执意要求男友开车送自己回家。男友无奈照办，不慎撞死人。男友构成交通肇事罪，女友也构成交通肇事罪。理由是，第一，客观上，女友指使醉酒者驾驶，由此给道路公共安全制造危险。这种指使行为是一种违法行为。并且，若没有这种违法行为，则男友不会导致死亡结果。第二，主观上，女友意识到醉酒者驾驶有可能撞死人，不过女友对这种死亡结果的发生不是故意心理，但是存在过失心理，属于过于自信的过失。

[注意］领导、女友虽然有指使行为，但不构成交通肇事罪的教唆犯或间接正犯。这是因为，教唆犯、间接正犯都是故意犯罪，教唆犯、间接正犯自身构成的罪名应是故意犯罪。例如，有故意杀人罪的教唆犯、危险驾驶罪（这是故意犯罪）的教唆犯，但没有过失致人死亡罪的教唆犯。又如，医生将一枚毒针交给护士，向护士隐瞒实情，让其打给一个病人，护士疏于查看，打给病人致其死亡。护士构成医疗事故罪，医生构成间接正犯，但不是医疗事故罪的间接正犯，而是故意杀人罪的间接正犯。

2.指使逃逸

例如，甲开车不慎撞死乙，重伤丙。甲欲救丙，路人丁指使甲逃逸别救人，甲逃逸，丙因未得到救助而死亡。甲构成交通肇事罪，属于“因逃逸致人死亡”。对丁如何处理？

% Source PDF physical page 325; see audit/source_coverage.json.
首先，甲故意不救人，属于故意的不作为犯罪。根据等价性的程度，甲构成遗弃罪，不过按照“因逃逸致人死亡”处理。丁指使甲不救人，构成遗弃罪的教唆犯。其次，指使犯罪分子逃逸，构成窝藏罪。\footnote{对窝藏罪的犯罪行为要扩大理解，不限于将犯罪分子藏在家里。参见窝藏罪一节内容。}最终，丁同时构成遗弃罪的教唆犯和窝藏罪，想象竞合，择一重罪论处。这是按照正常原理的分析。

[注意]司法解释的特殊规定：交通肇事后，\CJKunderwave{单位主管人员、车辆所有人、承包人或}\CJKunderwave{者乘车人}指使肇事人逃逸，致使被害人因得不到救助而死亡的，以交通肇事罪的共犯论处。注意该案件事实是，司机先已经构成交通肇事罪，之后四种指使者指使司机逃逸，因逃逸致人死亡，指使者构成司机的交通肇事罪的共犯。

一般犯罪结束后，有人参与进来，不可能构成该犯罪的承继共犯，只可能构成三种妨害司法犯罪（窝藏、包庇罪，掩饰、隐瞒犯罪所得罪，帮助伪造、毁灭证据罪）。当交通肇事罪成立后、结束后，有人参与进来，不可能构成交通肇事罪的承继共犯。所以，上述司法解释不符合基本原理，只能将其视为特殊规定，不能推广适用。注意：考试若考查司法解释的情形（四种指使者），就按该司法解释答题；若考查其他情形（其他指使者如路人），则按正常原理答题。（2006年第11题）

例如（2014年第13题），甲乘坐儿子乙（15周岁）的车，乙不慎撞死一人，重伤丙。甲指使乙逃逸，不要救丙。乙照办。丙因未得到及时救助而死亡。考试问，依照司法解释该如何处理？依照司法解释，乙作无罪处理，甲以交通肇事罪的共犯论处。考试又问，依照正常原理该如何处理？依照正常原理，在制造违法事实阶段，乙构成交通肇事罪，并且属于“因逃逸致人死亡”；在排除责任事由阶段，乙未满16周岁，排除责任。在制造违法事实阶段，就乙的“因逃逸致人死亡”而言，甲构成遗弃罪的教唆犯，甲同时构成窝藏罪；在排除责任事由阶段，甲符合条件。最终，对甲的两罪想象竞合，择一重罪论处。

\subsection*{二、危险驾驶罪}

第133条之一在道路上驾驶机动车，有下列情形之一的，处拘役，并处罚金：

\subsubsection*{（一）追逐竞驶，情节恶劣的；}

\subsubsection*{（二）醉酒驾驶机动车的；}

\subsubsection*{（三）从事校车业务或者旅客运输，严重超过额定乘员载客，或者严重超过规定时速}

行驶的；

\subsubsection*{（四）违反危险化学品安全管理规定运输危险化学品，危及公共安全的。}

机动车所有人、管理人对前款第三项、第四项行为负有直接责任的，依照前款的规定处罚。

有前两款行为，同时构成其他犯罪的，依照处罚较重的规定定罪处罚。[本条经《刑法修正案（九）》修正]

\subsubsection*{（一） 道路}

本罪的道路不限于公路、城市道路。（1）对某单位（如校园、居民小区等）管辖范围内的路段是否认定为“道路”，应当以其是否具有“公共性”，是否“允许社会机动车通行”作为判断标准。只允许单位内部机动车、特定来访机动车通行的，不认定为“道

% Source PDF physical page 326; see audit/source_coverage.json.
路”。（2）道路包括广场、公共停车场等用于公众通行的场所。（2015年第13题）

\subsubsection*{(二）醉酒驾车\protect\footnotemark}\footnotetext{2023年12月28日最高人民法院、最高人民检察院、公安部、司法部《关于办理醉酒危险驾驶刑事案件的意见》。}

1.第一种情形：80毫克的醉驾+严重情节，有罪。

醉驾是比酒驾更严重的程度，要求血液酒精含量达到80毫克/100毫升。

常考的严重情节有：（1）造成交通事故后逃逸的；（2）严重超员、超载、超速驾驶的；（3）毒驾；（4）从事客运且载有乘客的；（5）在高速公路上驾驶的。\footnote{其他严重情节有：（1）造成交通事故且负事故主要责任的；（2）驾驶重型载货汽车的；（3）无证驾驶的；（4）从事校车业务且载有师生的；（5）运输危险化学品、危险货物的；（6）逃避、阻碍公安机关依法检查的；（7）二年内曾因饮酒后驾驶机动车被查获或者受过行政处罚的。}

2.第二种情形：80毫克一150毫克的醉驾，无严重情节，无罪。

3.第三种情形：150毫克及以上的醉驾，无严重情节，有罪。

4.第四种情形：80毫克的醉驾+较轻情节，无罪。

较轻情节有：（1）出于急救伤病人员等紧急情况驾驶机动车，且不构成紧急避险的（例如，能够很快叫到网约车或救护车，却自己醉驾送医）；（2）在居民小区、停车场等场所因挪车、停车入位等短距离驾驶机动车的；（3）由他人驾驶至居民小区、停车场等场所短距离接替驾驶停放机动车的，或者为了交由他人驾驶，自居民小区、停车场等场所短距离驶出的。

[提示】醉酒后出于急救伤病人员等紧急情况，不得已驾驶机动车，构成紧急避险。

例如，甲的母亲心脏病突发，甲知道自己喝了酒，但由于情况紧急，不能很快叫到网约车或救护车，便驾车送母亲去医院。甲不构成危险驾驶罪，属于紧急避险。（2022年试题）

\subsubsection*{（三）主观要件}

本罪是故意犯罪。第一，主观上，只要求认识到自己处在酒精麻痹的状态，并不要求精确认识到自己血液中的酒精含量达到80mg/100ml或150mg/100ml。

第二，主观有醉酒驾驶机动车的意图。例如，刘某晚上跟朋友聚会喝醉酒。聚会结束，刘某坐在车里，打电话让老婆来接自己，打算让老婆开车送自已回家。在等老婆的时候，刘某坐在车里觉得太冷，便将发动机打开取暖。刘某没有驾驶汽车的意图，不构成危险驾驶罪。（2022年试题）

\subsubsection*{（四）共犯问题}

例1（教唆犯），甲喝醉酒，女友仍要求甲开车送自己回家。甲构成危险驾驶罪，女友是教唆犯。

例2（间接正犯），甲明知乙即将开车，暗中在其饮料中掺入酒精，乙没有意识到自已醉酒。甲构成危险驾驶罪的间接正犯，乙无罪。这也表明，醉酒驾驶不是真正身份犯。

例3（帮助犯），赵某和林某一起吃饭，林某喝了酒，赵某未喝酒。赵某开车送林某回家，途中林某执意要开车，于是赵某便停车，双方交换位置，林某醉酒驾车。赵某构成林某的危险驾驶罪的帮助犯。（2022年试题）

% Source PDF physical page 327; see audit/source_coverage.json.
[总结]危险驾驶行为的罪数关系\footnote{参见张明楷：《危险驾驶罪的基本问题》，载《政法论坛》2012年第6期。}

\InsertDiagram{figures/manual_table_p327_1.tex}

\subsection*{三、妨害安全驾驶罪}

第133条之二对\CJKunderwave{行驶中}的\CJKunderwave{公共交通工具}的驾驶人员\CJKunderwave{使用暴力或者抢控驾驶操纵装}\CJKunderwave{置}，干扰公共交通工具正常行驶，危及公共安全的，处一年以下有期徒刑、拘役或者管制，并处或者单处罚金。

% Source PDF physical page 328; see audit/source_coverage.json.
前款规定的驾驶人员在行驶的公共交通工具上擅离职守，与他人互殴或者殴打他人，危及公共安全的，依照前款的规定处罚。

有前两款行为，同时构成其他犯罪的，依照处罚较重的规定定罪处罚。[《刑法修正案

\subsubsection*{（十一）》增设]}

1.背景条件：行驶中的公共交通工具。

2.行为主体与行为方式：（1）乘客打司机、抢方向盘。（2）司机打乘客。（2022年试题）

3.总结罪数关系：（1）本罪是轻罪，危害公共安全的程度较小。如果危害公共安全的程度达到与放火罪、爆炸罪相当的程度，则构成以危险方法危害公共安全罪。

例1，甲抢夺方向盘，持续时间较长，导致公交车在大街上左摇右晃，行驶距离较长，严重危害道路公共安全。甲同时构成本罪和以危险方法危害公共安全罪，想象竞合，择一重罪论处。

例2，甲抢夺方向盘，想与全车人同归于尽，导致车掉进河里，死亡三人。甲构成本罪、以危险方法危害公共安全罪、故意杀人罪，想象竞合，择一重罪论处。（2024年试题）

（2）本罪可以与劫持汽车罪想象竞合。例如，甲抢夺方向盘，劫持公交车，驶向自己的目的地。甲构成本罪、劫持汽车罪，想象竞合，择一重罪论处，定劫持汽车罪。

（3）本罪可以与交通肇事罪想象竞合。例如，甲抢夺方向盘，短时间内导致公交车失控，轧死一名行人。甲构成本罪和交通肇事罪，想象竞合，择一重罪论处，定交通肇事罪。

[总结]驾驶中的罪名区分

\InsertDiagram{figures/manual_table_p328_1.tex}

\section{恐怖型犯罪}

\subsection*{一、组织、领导、参加恐怖组织罪（第120条）}

1.恐怖组织与黑社会性质组织的区别：前者有政治目的，后者没有政治目的。

2.成立本罪只要求组织、领导、参加恐怖组织，不要求实施恐怖活动。

% Source PDF physical page 329; see audit/source_coverage.json.
3.犯本罪又实施其他犯罪的，数罪并罚。

\subsection*{二、帮助恐怖活动罪（第 120 条之一）}

1.资助，仅限于物质性资助，不包括精神性资助。

2.帮助行为被正犯化。本罪是将帮助行为作为实行行为对待，将帮助犯作为正犯对待，成为独立罪名。这种规定产生两个法律后果：（2016年第56题）

（1）在定罪条件上，不需要遵守共犯从属性，也即，被帮助的人即使不构成犯罪，帮助行为也构成犯罪，因为帮助行为自身已经是实行行为。

（2）在共犯关系上，由于原帮助行为成为实行行为，那么教唆它、帮助它的行为，就可以成立教唆犯、帮助犯。例如，甲教唆乙资助恐怖组织，乙答应照办。乙构成帮助恐怖活动罪，甲构成该罪的教唆犯。

3.既遂。本罪的既遂标准不是所资助的恐怖组织实施了恐怖活动，而是提供的资助被恐怖组织接收。例如，甲为恐怖组织招募了人员，恐怖组织接收了该人员，本罪就既遂。

4.总结。与本罪类似的还有资助危害国家安全犯罪活动罪（第107条），也是帮助行为正犯化。例如（2016年第53题），丁资助了林某从事危害国家安全犯罪活动，林某尚未实施便被抓。丁构成资助危害国家安全犯罪活动罪的既遂，而非未遂。

\subsection*{三、准备实施恐怖活动罪}

第120条之二有下列情形之一的，处五年以下有期徒刑、拘役、管制或者剥夺政治权利，并处罚金；情节严重的，处五年以上有期徒刑，并处罚金或者没收财产：

\subsubsection*{（一）为实施恐怖活动准备凶器、危险物品或者其他工具的；}

\subsubsection*{（二）组织恐怖活动培训或者积极参加恐怖活动培训的；}

\subsubsection*{（三）为实施恐怖活动与境外恐怖活动组织或者人员联络的；}

\subsubsection*{（四）为实施恐怖活动进行策划或者其他准备的。}

有前款行为，同时构成其他犯罪的，依照处罚较重的规定定罪处罚。[本条为《刑法修正案（九）》增设]

本罪的特点是，将预备行为正犯化，成为一种独立的预备罪。

1.教唆、帮助实施本罪，构成本罪的教唆犯、帮助犯。当然，以被教唆、帮助者实施了本罪为前提。

2.为实施本罪而进行的预备行为，如果具有现实的危险（该危险只要求抽象缓和，不要求紧迫直接），则构成本罪的预备犯。

例1，为了组织恐怖活动培训，已经联系了讲师，准备了培训场地，此时被抓捕，构成本罪的犯罪预备，成立预备犯。

例2，为了准备危险物品，阅读相关书籍或在网上查阅相关资料，不具有现实危险，不构成本罪的预备犯。

例3，为了购买凶器而打工挣钱，不具有现实危险，不构成本罪的预备犯。

% Source PDF physical page 330; see audit/source_coverage.json.
\section{枪支类犯罪}

\subsection*{一、枪支的认定}

1.行政法将枪支的认定标准设定的很低，不要求较大杀伤力，据此，天津大妈赵春华案中的游戏枪也是枪支。该案之后，刑法的司法解释规定，刑法上的枪支要求具有较大杀伤力。\footnote{2018年3月30日最高人民法院、最高人民检察院《关于涉以压缩气体为动力的枪支、气枪铅弹刑事案件定罪量刑问题的批复》。}此外，弓弩不是枪支。

2.司法解释规定，就本节的枪支类犯罪而言，以火药为动力发射弹药的大口径武器，可认定为枪支。基于此，土炮、大炮也属于枪支。这是一种扩大解释。

\subsection*{二、盗窃、抢夺枪支、弹药、爆炸物、危险物质罪（第127条第1款）}

1.本罪是危险犯，成立本罪不要求造成严重后果。但是不意味着一旦实施就既遂，既遂标准是实际取得枪支、弹药、爆炸物、危险物质。

2.骗取行为，不构成本罪，可成立诈骗罪。

3.认识错误：行为人以普通盗窃的故意，盗窃了枪支，定普通盗窃罪既遂，但如果事后持有，就构成非法持有枪支罪，数罪并罚。（2007年第61题）

\subsection*{三、非法持有、私藏枪支、弹药罪（第128条第1款）}

1.非法持有：是指没有合法根据而实际占有或控制枪支、弹药。非法替他人保管枪支、弹药，也属于非法持有。

2.非法私藏：是指依法配备、配置枪支、弹药的人员，在配备、配置枪支、弹药的条件消除后，违反枪支管理法律、法规的规定，私自藏匿所配备、配置的枪支、弹药且拒不交出的行为。

3.本罪与非法储存枪支、弹药罪，非法携带枪支、弹药危及公共安全罪的区分总结：

（1）非法储存与非法持有的区分：数量大小。储存是大量，持有是少量。

（2）非法携带，是指将枪支置于现实控制之下，进入公共场所或公共交通工具，危及公共安全。行为人既可以是不符合配备、配置枪支条件的人，也可以是符合配备、配置枪支条件的人。

\subsection*{四、非法出租、出借枪支罪}

第128条第2款依法配备公务用枪的人员，非法出租、出借枪支的，依照前款的规定处罚。

第128条第3款依法配置枪支的人员，非法出租、出借枪支，造成严重后果的，依照第一款的规定处罚。

1.行为主体是特殊主体：（1）依法配备公务用枪的人员和单位；（2）依法配置枪支的人员和单位。（2012年第58题）

% Source PDF physical page 331; see audit/source_coverage.json.
2.行为方式

（1）出租：是指在一段期限内有偿提供给他人使用。如果是永久性地有偿转让给他人，则属于非法买卖枪支罪。

（2）出借：一般是指在一段期限内无偿提供给他人使用。但是，非法将枪支赠与他人，比非法出借的危害性更大，更值得处罚。因此，这里的“出借”应作扩大解释，包括永久性地无偿提供给他人使用，也即包括赠与。

[注意］根据司法解释，依法配备公务用枪的人员，将公务用枪用作借债质押物，使枪支处于非法持枪人的控制、使用之下，严重危害公共安全，以非法出借枪支罪论处；对接受者，构成犯罪的，以非法持有枪支罪论处。

\subsection*{五、丢失枪支不报罪}

第129条依法配备公务用枪的人员，丢失枪支不及时报告，造成严重后果的，处三年以下有期徒刑或者拘役。

1.行为主体是依法配备公务用枪的人员。猎人不属于本罪主体。（2017年第12题）

2.丢失枪支的“丢失”应作扩大解释，指非自愿失去对枪支的占有控制，包括遗失、被盗、被抢。

3.本罪的行为是不报告，是真正不作为犯。不作为犯存在共犯现象。例如，警察甲与乙外出追捕逃犯。其间，甲的枪被逃犯抢去后，意欲及时报告。乙唆使甲不报告：“别报告了，不嫌丢人，以后需要用枪，用兄弟我的。”于是甲没报告。后来该枪落入不法分子之手，造成了严重后果。乙构成教唆犯，甲构成实行犯。

4.主观：故意，但对造成严重后果，不要求行为人有认识，只要求有认识可能性。这种“严重后果”在理论上被称为客观超过要素或客观处罚条件。

\subsection*{六、非法制造、买卖、运输、邮寄、储存\CJKunderwave{枪支、弹药、爆炸物}罪；非法制造、买卖、运输、储存危险物质罪}

第125条第1、2款非法制造、买卖、运输、邮寄、储存枪支、弹药、爆炸物的，处三年以上十年以下有期徒刑；情节严重的，处十年以上有期徒刑、无期徒刑或者死刑。

非法制造、买卖、运输、储存毒害性、放射性、传染病病原体等物质，危害公共安全的，依照前款的规定处罚。

1．“买卖”。第一，既包括卖，也包括买。买，既包括为卖而买，也包括为自用而买。第二，买卖的方式，包括钱物交易，也包括物物交换。例如，甲用枪，乙用毒品，相互交换。甲构成非法买卖枪支罪，乙构成贩卖毒品罪和非法买卖枪支罪，数罪并罚。（2019年试题）

2.本罪是选择性罪名，可以分拆使用。如果产生认识错误，看客观。例如，想买枪，却买到子弹，定非法买卖弹药罪。（2020年试题）

\section{安全生产领域犯罪}

\subsection*{一、重大责任事故罪与重大劳动安全事故罪}

二者的区分标准：造成事故的原因不同，前者是在生产、作业过程中违反安全管理规定（操作、作业有问题）；后者是安全生产条件不符合规定（设备、条件有问题）。

% Source PDF physical page 332; see audit/source_coverage.json.
\subsection*{二、危险作业罪（第134条之一）}

该罪是指，在生产、作业中，违反有关安全管理的规定，具有发生重大伤亡事故的现实危险的行为。

1.客观要件：要求造成具体危险。本罪是具体危险犯（2021年试题）。若造成实害结果，则构成重大责任事故罪。

2.主观要件：本罪是故意犯罪。

\InsertDiagram{figures/manual_table_p332_1.tex}

\subsection*{三、强令、组织他人违章冒险作业罪}

第134条第2款强令他人违章冒险作业，或者明知存在重大事故隐患而不排除，仍冒险组织作业，因而发生重大伤亡事故或者造成其他严重后果的，处五年以下有期徒刑或者拘役；…[《刑法修正案（十一）》修订]

[注意】本罪是过失犯罪，要求造成实害结果。

\subsection*{四、不报、谎报安全事故罪（第139条之一）}

该罪是指，在安全事故发生后，负有报告职责的人员不报或者谎报事故情况，贻误事故抢救，情节严重的行为。

成立本罪，要求不报、谎报行为导致贻误事故抢救。这表明：

第一，如果行为人尽力抢救了，没有贻误抢救，则即使不报、谎报，也不构成本罪。（2025年试题）。

第二，如果抢救工作全部结束后，行为人不报、谎报，不构成本罪。（2024年试题）
